\documentclass[12pt, letterpaper]{article}

% --- Paquetes básicos ---
\usepackage[spanish, es-tabla]{babel}
\usepackage[utf8]{inputenc}
\usepackage[T1]{fontenc}
\usepackage{newpxtext,newpxmath} % Fuente estilo Palatino
\usepackage[margin=3cm]{geometry}
\usepackage{titlesec}
\usepackage[autostyle]{csquotes}

% --- Para dibujar la V de Gowin ---
\usepackage{tikz}
\usetikzlibrary{positioning,babel,calc}

% --- Configuración de la sección (1. IDEA DE...) ---
\titleformat{\section}
  {\normalfont\large\scshape\centering} 
  {\thesection.}
  {0.5em}
  {}

\begin{document}

% --- Encabezado Manual ---
\begin{center}
    \vspace{0.8cm}
    {\Large \bfseries BENFORD'S LAW: THEORY AND APPLICATIONS} \\
    \vspace{0.4cm}
\end{center}

\vspace{1.5cm}

\noindent
\textbf{Nivel descriptivo} \newline

\noindent
Título: Benford's Law: Theory and Applications \newline

\noindent
Autor: Steven J. Miller \newline

\noindent
Año: 2015 \newline

\noindent
Nombre del tema: Estudio matemático y aplicado de la Ley de Benford. \newline

\noindent
Clasificación: Teórica y aplicada; el libro desarrolla los fundamentos de la Ley de Benford y muestra cómo puede usarse para detectar patrones no uniformes en datos reales.\newline

\noindent
Resumen en una oración: El libro explica por qué muchos conjuntos de datos naturales no distribuyen sus dígitos iniciales de forma uniforme. \newline

\noindent
Argumento central: Miller sostiene que la Ley de Benford aparece con frecuencia en datos reales y que puede ser usada como herramienta para analizar la estructura de ciertos conjuntos numéricos.\newline

\noindent
Problemas con el argumento: La ley no funciona en todos los casos, así que su aplicación depende mucho del tipo de datos estudiados. Si el conjunto no tiene una estructura adecuada, la distribución benfordiana puede no ser representativa.
\newpage 

\noindent
\textbf{Nivel analítico} \newline

\noindent
Conexión con mi proyecto: Aunque no trata directamente sobre contraseñas, este texto resulta útil porque introduce una forma de pensar en distribuciones no uniformes que podría inspirar el análisis de patrones dentro de secuencias de contraseñas.\newline

\noindent
Lo que NO dice el autor: El libro no estudia seguridad informática, ni analiza cómo se comportan las contraseñas humanas, ni ofrece métodos para medir entropía de contraseñas. \newline

\noindent
Contraste con otra fuente: Frente a trabajos centrados en la predicción de contraseñas, aquí el interés está en una ley estadística general que puede aplicarse a diversos tipos de datos, no a un problema de ciberseguridad en particular.   \newline 

\noindent
Evaluación de aplicabilidad: 2. Es una fuente útil como apoyo conceptual, pero su relación con el proyecto es indirecta y requiere una adaptación considerable. \newline

\noindent
Preguntas que le haría al autor: ¿Cree que la lógica estadística de la Ley de Benford podría servir como inspiración para analizar frecuencias de caracteres o posiciones dentro de contraseñas? \newline

\noindent
Resumen argumentativo: El libro intenta mostrar que muchos datos reales siguen patrones numéricos sorprendentes y no uniformes. Para el proyecto, su aporte principal no está en el contenido específico, sino en la idea de que una distribución empírica puede revelar regularidades útiles para interpretar datos complejos.

\end{document}
